```latex
% ch4_implementation.tex

\chapter{Implémentation}
\label{ch:implementation}

Ce chapitre présente l'implémentation distribuée de StreamKM++ avec Apache Spark. La logique algorithmique décrite au chapitre~\ref{ch:algorithmes} reste indépendante de Spark et est regroupée dans les packages \texttt{streamkm.core} et \texttt{streamkm.coreset}. La couche \texttt{streamkm.spark} prend en charge la distribution des données, la construction des coresets locaux, leur fusion et l'évaluation du clustering obtenu.

L'implémentation repose sur deux niveaux de traitement. Chaque partition construit d'abord un résumé local des données qui lui sont attribuées. Ces résumés sont ensuite fusionnés afin d'obtenir un coreset global sur lequel le clustering final est calculé. Structured Streaming reprend ce même principe pour traiter successivement plusieurs micro-batches.

% ---------------------------------------------------------------------------
\section{Construction distribuée du coreset}
\label{sec:impl-dataflow}
% ---------------------------------------------------------------------------

\subsection{Construction d'un coreset par partition}

La première étape consiste à répartir la construction du coreset entre les partitions du RDD. La méthode \texttt{StreamKMPlusPlus.partialCoresets} utilise \texttt{mapPartitionsWithIndex} afin que chaque partition construise indépendamment son propre résumé.

\begin{verbatim}
def partialCoresets(
    rdd: RDD[Point],
    p: StreamKMParams
): RDD[Array[Point]] = {

  rdd.mapPartitionsWithIndex { (idx, it) =>
    val bm = new BucketManager(p.m, p.seed + idx)
    bm.insertAll(it)
    Iterator.single(bm.extractCoreset())
  }
}
\end{verbatim}

Chaque partition instancie un \texttt{BucketManager}, lui transmet les observations qu'elle contient puis extrait le coreset obtenu. La graine utilisée est dérivée de l'indice de la partition afin de disposer d'une construction locale reproductible pour un partitionnement donné.

Cette organisation correspond naturellement au modèle distribué de Spark. Les partitions peuvent être traitées indépendamment avant la phase de fusion, ce qui évite de centraliser la construction du coreset sur le driver.

\subsection{Fusion des résumés locaux}

Les coresets construits sur les différentes partitions doivent ensuite être combinés afin de produire un résumé global. La méthode \texttt{mergeCoresets} applique pour cela une réduction arborescente.

\begin{verbatim}
def mergeCoresets(
    rdd: RDD[Point],
    p: StreamKMParams,
    depth: Int = 2
): Array[Point] = {

  val perPartition: RDD[BucketManager] =
    rdd.mapPartitionsWithIndex { (idx, it) =>
      val bm = new BucketManager(p.m, p.seed + idx)
      bm.insertAll(it)
      Iterator.single(bm)
    }

  perPartition
    .treeReduce((a, b) => a.merge(b), depth)
    .extractCoreset()
}
\end{verbatim}

Chaque partition produit ici un \texttt{BucketManager} contenant son résumé local. Les structures obtenues sont ensuite combinées à l'aide de \texttt{treeReduce}. L'opération de fusion repose directement sur la méthode \texttt{BucketManager.merge}, déjà utilisée par la stratégie \emph{merge-and-reduce} décrite au chapitre~\ref{ch:algorithmes}.

Cette approche permet de réutiliser la même logique de fusion indépendamment du niveau auquel elle est appliquée. La couche Spark organise ainsi l'exécution distribuée sans modifier les mécanismes internes de construction du coreset.

Par rapport à une récupération de l'ensemble des résumés sur le driver suivie d'une fusion séquentielle, la réduction arborescente répartit une partie des opérations de combinaison entre plusieurs niveaux.

% ---------------------------------------------------------------------------
\section{Traitement avec Structured Streaming}
% ---------------------------------------------------------------------------

La méthode \texttt{StreamKMPlusPlus.run} étend la construction distribuée au traitement d'un flux avec Structured Streaming. Un \texttt{BucketManager} global est maintenu sur le driver afin de conserver un résumé des micro-batches déjà traités.

\begin{verbatim}
def run(
    source: DataFrame,
    p: StreamKMParams,
    onModel: KMeansModel => Unit,
    featuresCol: String = "features"
): StreamingQuery = {

  val global = new BucketManager(p.m, p.seed)

  source.writeStream.foreachBatch { (batch, batchId) =>

    val points = batch.rdd.map { row =>
      Point(
        row.getAs[Seq[Double]](featuresCol).toArray
      )
    }

    val partial =
      mergeCoresets(
        points,
        p.copy(seed = p.seed + batchId)
      )

    partial.foreach(global.insert)

    val model =
      KMeans.fit(
        global.extractCoreset(),
        p.k,
        p.nRestarts,
        seed = p.seed
      )

    onModel(model)

  }.start()
}
\end{verbatim}

Chaque micro-batch est d'abord converti en RDD de \texttt{Point}. Un coreset est ensuite construit de manière distribuée à partir des partitions du lot. Le résumé obtenu est intégré au \texttt{BucketManager} global, qui représente progressivement l'ensemble des observations déjà traitées.

Le clustering est calculé à partir du coreset courant après l'intégration du micro-batch. La méthode \texttt{onModel} permet ensuite d'exploiter le modèle produit sans imposer un mécanisme de sortie particulier.

Les observations brutes sont considérées comme ayant un poids unitaire. Les poids supérieurs à un apparaissent uniquement lors de la construction des coresets et correspondent à la masse représentée par chaque point du résumé.

% ---------------------------------------------------------------------------
\section{Correspondance avec l'implémentation Flink}
% ---------------------------------------------------------------------------

L'implémentation Spark conserve les principales étapes de l'approche distribuée proposée avec Flink par Bitsakis \cite{bitsakis2018}, mais leur réalisation est adaptée au modèle d'exécution de Spark.

\begin{table}[ht]
\centering
\small
\begin{tabular}{p{3.5cm}p{4.5cm}p{4.8cm}}
\toprule
\textbf{Opérateur Flink} &
\textbf{Rôle} &
\textbf{Implémentation Spark} \\
\midrule

\texttt{HdfsSource} &
Lecture distribuée des données &
Source Spark partitionnée \\

\texttt{HandleWatermark} &
Identification d'une frontière de traitement &
Frontière fournie par les micro-batches de Structured Streaming \\

\texttt{FlatMapToPoint} &
Conversion des données en points &
Transformation en objets \texttt{Point} \\

\texttt{FlatMapToPartialCoresets} &
Construction locale du coreset &
\texttt{mapPartitionsWithIndex} \\

\texttt{FlatMapToFinalCoreset} &
Fusion des coresets partiels &
\texttt{treeReduce} et \texttt{BucketManager.merge} \\

\texttt{FlatMapToKmeansPP} &
Calcul du clustering final &
\texttt{KMeans.fit} sur le coreset global \\

\texttt{HdfsSink} &
Production des centres &
Fonction de sortie \texttt{onModel} \\

\bottomrule
\end{tabular}
\caption{Correspondance entre les principales étapes des implémentations Flink et Spark.}
\label{tab:flink-spark}
\end{table}

La principale différence concerne la fusion des résumés intermédiaires. Dans l'implémentation Spark, les coresets locaux peuvent être combinés au moyen d'une réduction arborescente. Le découpage en micro-batches fournit également une frontière explicite entre les différents lots traités.

% ---------------------------------------------------------------------------
\section{Évaluation distribuée du coût}
\label{sec:impl-evaluation}
% ---------------------------------------------------------------------------

La qualité du clustering doit être évaluée sur l'ensemble des données originales et non uniquement sur le coreset. Le coût utilisé dans les expériences correspond à la somme des distances euclidiennes au carré entre chaque observation et son centre le plus proche.

Cette opération est réalisée de manière distribuée par la méthode \texttt{CostEvaluator.cost} :

\begin{verbatim}
def cost(
    points: RDD[Point],
    centers: Array[Array[Double]],
    depth: Int = 2
): Double = {

  val bcCenters =
    points.sparkContext.broadcast(centers)

  points
    .mapPartitions { it =>
      Iterator.single(
        KMeans.cost(
          it.toArray,
          bcCenters.value
        )
      )
    }
    .treeReduce(_ + _, depth)
}
\end{verbatim}

Les centres sont d'abord diffusés aux exécuteurs à l'aide d'une variable \texttt{broadcast}. Chaque partition calcule ensuite localement la contribution de ses observations au coût total. Les coûts partiels sont enfin additionnés à l'aide de \texttt{treeReduce}.

Cette méthode évite de rapatrier les observations sur le driver pour effectuer l'évaluation. Elle est utilisée dans les expériences du chapitre~\ref{ch:evaluation} afin de comparer les différentes configurations sur la même métrique et sur l'ensemble du jeu de données.

% ---------------------------------------------------------------------------
\section{Périmètre de l'implémentation}
% ---------------------------------------------------------------------------

La thèse de Bitsakis décrit également une variante permettant de produire périodiquement plusieurs clusterings pour différentes valeurs de \(k\), ainsi qu'un mécanisme d'interrogation de l'état du traitement.

Ces fonctionnalités n'ont pas été retenues dans le périmètre de l'implémentation réalisée. Le travail se concentre sur la construction distribuée de StreamKM++, son adaptation à Structured Streaming et l'évaluation expérimentale de la qualité et des performances du portage vers Spark.

Structured Streaming permet néanmoins de recalculer un modèle après chaque micro-batch à partir du coreset courant, comme le montre la méthode \texttt{run}. Cette fonctionnalité fournit le mécanisme nécessaire à la production progressive de résultats sans ajouter de gestion spécifique de requêtes périodiques.

% ---------------------------------------------------------------------------
\section{Synthèse du portage de Flink vers Spark}
\label{sec:portage-flink-spark}
% ---------------------------------------------------------------------------

Le portage vers Spark conserve les principes algorithmiques de l'implémentation distribuée de StreamKM++, tout en adaptant leur orchestration au modèle d'exécution du framework.

La construction locale des coresets est naturellement associée aux partitions Spark. Les résumés produits peuvent ensuite être combinés à l'aide de \texttt{treeReduce}, en réutilisant directement les opérations de fusion définies dans \texttt{BucketManager}. Structured Streaming permet quant à lui de traiter successivement les différents micro-batches et de maintenir un coreset représentant les données déjà consommées.

L'évaluation du coût suit la même logique distribuée : les centres sont diffusés aux partitions et les contributions locales sont agrégées pour produire la somme des erreurs quadratiques sur l'ensemble des observations.

Cette architecture permet ainsi de conserver une séparation claire entre la logique algorithmique de StreamKM++ et les mécanismes de distribution propres à Spark. Les expérimentations présentées dans le chapitre suivant permettent d'évaluer la qualité des coresets obtenus ainsi que le comportement de cette implémentation lorsque la taille des données, le partitionnement et les ressources de calcul varient.
```
